\metatitle{M-convexity and completion polynomials}
\metadescription{The general M-convexity theorem for one-color count vectors,
completion polynomials, their Lorentzian conjecture, and block factorization.}

\section[rank]{M-convexity and completion polynomials}

For every colored matroid, counting the positions of one color in parts
of a partition gives an M-convex set of vectors. We prove this general
result, then study the polynomials that count how many words give each
vector.

Let $\mathcal F$ be a colored matroid on $E$ with colors $1,\dotsc,k$ and
quotas $r_1,\dotsc,r_k$. We write
\[
 B(u,c)\coloneqq\{e\in E:u_e=c\}.
\]
The \href{intro.htm}{color-projection proposition} shows that
$\{B(u,c):u\in\mathcal F\}$ is the basis family of an ordinary matroid,
which we denote by $M_c$.
\label{mergingColors}
Retaining color $c$ and merging all the others gives the two-color
matroid associated with $M_c$ and its dual. More general
\href{operations.htm#colorMerging}{color merging can destroy exchange}.

\section[mConvexCounts]{M-convex count vectors}

Fix a color $c$ and a partition $E=D_1\sqcup\dotsb\sqcup D_m$.
We record how many positions of each part receive color $c$:
\[
 \defin{\mathcal V_c}\coloneqq
 \{(|B(u,c)\cap D_i|)_{i=1}^m:u\in\mathcal F\}.
\]

\begin{definition}[M-convexity]
A nonempty finite set $\mathcal V\subseteq\mathbb Z_{\geq0}^m$ of
constant coordinate sum is \defin{M-convex} if, whenever
$\mathbf a,\mathbf b\in\mathcal V$ and $a_i>b_i$, there is an index $j$
with $a_j<b_j$ such that
\[
 \begin{gathered}
 \mathbf a-\mathbf e_i+\mathbf e_j\in\mathcal V,\\
 \mathbf b+\mathbf e_i-\mathbf e_j\in\mathcal V.
 \end{gathered}
\]
Here $\mathbf e_i$ is the $i$th unit vector.
\end{definition}

This is symmetric basis exchange for integer vectors rather than just
zero-one vectors \cite{Murota2003}.

\begin{theorem}[M-convex count vectors]
The set $\mathcal V_c$ is M-convex. More precisely, put
\[
 f(J)\coloneqq\operatorname{rk}_{M_c}
 \left(\bigcup_{i\in J}D_i\right)
 \qquad(J\subseteq[m]).
\]
The set $\mathcal V_c$ consists exactly of the integer vectors
$\mathbf a\in\mathbb Z_{\geq0}^m$ satisfying
\[
 \begin{gathered}
 \sum_{i=1}^m a_i=r_c,\\
 \sum_{i\in J}a_i\leq f(J)\qquad(J\subseteq[m]).
 \end{gathered}
\]
\end{theorem}

\begin{proof}
We show that the displayed inequalities characterize the counts of bases
of $M_c$. Necessity follows from the definition of matroid rank.

For sufficiency, take an integer vector $\mathbf a$ satisfying them.
Let $N$ be the partition matroid in which an independent set contains at
most $a_i$ elements of $D_i$. The matroid intersection theorem
\cite{Oxley2011} gives a common independent set of size $r_c$ if,
for every $X\subseteq E$,
\[
 \operatorname{rk}_{M_c}(X)+
 \sum_i\min(a_i,|D_i\setminus X|)\geq r_c.
\]
To verify this, put $J=\{i:a_i>|D_i\setminus X|\}$. Then
\[
 \sum_{i\in J}a_i\leq f(J)
 \leq\operatorname{rk}_{M_c}(X)+\sum_{i\in J}|D_i\setminus X|.
\]
Adding $\sum_{i\notin J}a_i$ gives the required inequality.
A common independent set of size $r_c$ is a basis of $M_c$ and has
exactly $a_i$ elements in each $D_i$.

The function $f$ is integer-valued, nondecreasing, and submodular, with
$f(\varnothing)=0$. Hence the displayed count set is the basis set of
an integral polymatroid. The symmetric exchange theorem for discrete
polymatroids gives M-convexity; see
\href{https://arxiv.org/html/math/0307236v1}{J.~Herzog and T.~Hibi,
Theorems 3.4 and 4.1}.
\end{proof}

Taking singleton parts recovers ordinary symmetric basis exchange.
Taking $D_1=S$ and $D_2=E\setminus S$ shows that the possible values of
$|B(u,c)\cap S|$ form an integer interval.

\begin{example}[Counts and multiplicities]
Take all two-color words on four positions with quotas $(2,2)$, and
partition the positions as $\{1,2\}\sqcup\{3,4\}$. The color-$1$
count vectors are
\[
 \mathcal V_1=\{(2,0),(1,1),(0,2)\}.
\]
Their multiplicities are $1,4,1$: a word either puts both $1$'s in
one part or chooses one position from each part.
\end{example}

\section[completionPolynomials]{Completion polynomials}

The count-vector theorem describes which vectors occur. A completion
polynomial also records their multiplicities.

Fix a color $c$ and positive weights $w_{e,a}$ for the assignments with
$a\ne c$. Define the \defin{weighted completion polynomial}
\[
 \begin{gathered}
 \defin{Q_c(\mathbf x;\mathbf w)}\coloneqq\\
 \sum_{u\in\mathcal F}
 \left(\prod_{e:u_e\ne c}w_{e,u_e}\right)
 \prod_{e:u_e=c}x_e.
 \end{gathered}
\]
It is homogeneous of degree $r_c$. When all weights are $1$, the
coefficient of $\prod_{e\in B}x_e$ counts the words whose color-$c$
class is $B$.
For a polynomial with nonnegative coefficients, its \defin{exponent support}
is the set of exponent vectors of monomials with positive coefficient.

\begin{corollary}[M-convex support]
For every positive choice of weights, identifying $x_e=z_i$ for
$e\in D_i$ gives a polynomial whose exponent support is exactly
$\mathcal V_c$, and is therefore M-convex.
\end{corollary}

\begin{proof}
Every word contributes a positive coefficient to its count-vector
monomial, so the support consists of all and only the vectors in
$\mathcal V_c$.
\end{proof}

In the preceding example, the unweighted polynomial after identifying
variables within each part is $z_1^2+4z_1z_2+z_2^2$.

\section[lorentzianConjecture]{The Lorentzian conjecture}

\begin{definition}[Lorentzian polynomial]
Let $f$ be a nonzero homogeneous polynomial of degree $d$ with nonnegative
coefficients.
For $d\geq2$, we call $f$ \defin{Lorentzian} if its exponent support is
M-convex and every partial derivative of total order $d-2$ has a Hessian
with at most one positive eigenvalue. The Hessian is the matrix of second
partial derivatives. Nonnegative homogeneous polynomials of degree zero
or one are also Lorentzian \cite{BrandenHuh2020}.
\end{definition}

\begin{conjecture}[Lorentzian completion polynomials]
For every colored matroid $\mathcal F$, every color $c$, and every choice
of positive assignment weights, $Q_c(\mathbf x;\mathbf w)$ is Lorentzian.
\end{conjecture}

Its support is already M-convex by the color-projection proposition.
The conjecture asks for the stronger inequalities on completion counts.
The \href{reductions.htm}{reductions chapter} proves that its universal
weighted and unweighted forms are equivalent, and that binary-list
families with distinguished quota two suffice, even with a fixed number
of colors.
In two colors, each basis has a unique completion. If $a$ is the other
color, then
\[
 Q_c(\mathbf x;\mathbf w)=
 \left(\prod_{e\in E}w_{e,a}\right)
 \sum_{B\in\mathcal B(M_c)}\prod_{e\in B}\frac{x_e}{w_{e,a}},
\]
so the conjecture follows from the Lorentzian theorem for ordinary
matroid basis polynomials \cite{BrandenHuh2020}.

Nonnegative linear substitutions preserve Lorentzianity, so the conjecture
would also make the partition-count polynomials Lorentzian. With two parts
and all weights equal to $1$, write this polynomial as
\[
 \sum_{j=0}^{r_c}a_j z_1^jz_2^{r_c-j}.
\]
The resulting inequality is
\[
 \left(\frac{a_j}{\binom{r_c}{j}}\right)^2\geq
 \frac{a_{j-1}}{\binom{r_c}{j-1}}
 \frac{a_{j+1}}{\binom{r_c}{j+1}}
 \qquad(1\leq j<r_c).
\]
Thus the conjecture strengthens the integer-interval statement to
ultra-log-concavity. The \href{log-concavity.htm}{width-three poset chapter}
proves Lorentzianity for the complement of two disjoint matroid bases and
applies it to descent counts.

\section[blockFactorization]{Fixed block degrees force factorization}

We record an observation about Lorentzian polynomials that is useful when
choosing a refinement of a completion polynomial. Fixed degrees in separate
blocks of variables leave no room for interaction between those blocks.

Partition the variables into disjoint blocks $X_1,\dotsc,X_m$. We say that
$f$ has \defin{fixed block degrees} $(d_1,\dotsc,d_m)$ if every monomial
of $f$ has total degree $d_a$ in the variables of $X_a$, for every $a$.

\begin{theorem}[Block factorization]
Let $f\ne0$ be Lorentzian with fixed block degrees $(d_1,\dotsc,d_m)$.
Then
\[
 f=C\prod_{a=1}^m f_a(X_a),\qquad C>0,
\]
where each $f_a$ is a nonzero Lorentzian polynomial of degree $d_a$.
The factors are unique up to positive scalar multiples. Conversely, every
such product is Lorentzian. Degree-zero blocks contribute constants.
\end{theorem}

\begin{proof}
We first show that the support is a product, then factor the coefficients.
Let $\mathcal S$ be the exponent support and $\mathcal S_a$ its projection
to block $X_a$. An M-convex exchange from a surplus coordinate in $X_a$
must end in $X_a$, since an exchange between blocks would change their
fixed degrees. Repeated exchanges can therefore replace one block of a
supported vector by the same block of any other supported vector, keeping
all other blocks fixed. Consequently
\[
 \mathcal S=\mathcal S_1\times\dotsb\times\mathcal S_m.
\]
Each $\mathcal S_a$ is M-convex, and its graph of unit exchanges is connected.

Write the coefficients in factorial-normalized form,
\[
 f=\sum_{\alpha\in\mathcal S}a_\alpha\frac{x^\alpha}{\alpha!},
 \qquad a_\alpha>0.
\]
Choose adjacent vectors $p,p'$ in one projected support and $q,q'$ in
another, and fix the remaining blocks at a supported tuple $\eta$.
Differentiate by the coordinatewise minima of the adjacent pairs in
these two blocks, and by $\eta$ elsewhere. The derivative has bidegree
$(1,1)$ and hence Hessian
\[
 \begin{pmatrix}0&M\\M^{\mathsf T}&0\end{pmatrix}.
\]
This matrix has one positive eigenvalue for each nonzero singular value
of $M$. Lorentzianity therefore forces $\operatorname{rank}M\leq1$.
The four chosen terms give the identity
\[
 a_{p,q,\eta}a_{p',q',\eta}
 =a_{p',q,\eta}a_{p,q',\eta}.
\]
The entries are precisely the normalized coefficients: each surviving
term has exponent one in each of the two remaining blocks.

Thus a coefficient ratio along an exchange edge in one block is
independent of the other blocks. Following paths in their connected
exchange graphs gives
\[
 a_\alpha=C\prod_a h_a(\alpha|_{X_a})
\]
for positive functions $h_a$. Since $\alpha!$ also factors by blocks,
this proves the asserted polynomial factorization and its uniqueness.
Differentiating by a supported monomial in every other block leaves a
positive multiple of $f_a$, so each factor is Lorentzian. The converse
is closure under products. We use the support and derivative criteria
and product closure from \cite{BrandenHuh2020}.
\end{proof}

\begin{remark}[Proof status]
The argument above is a written proof. Independent review and formal
verification have not yet been completed.
\end{remark}

\subsection[stableBlockFactorization]{Stable polynomials}

A nonzero polynomial with complex coefficients is \defin{stable} if it
does not vanish when all variables have positive imaginary part. It is
\defin{real stable} when its coefficients are real.

\begin{corollary}[Stable block factorization]
If $f$ is stable with fixed block degrees, then
\[
 f=c\prod_a f_a(X_a),\qquad c\in\mathbb C\setminus\{0\},
\]
where the $f_a$ are real stable with nonnegative coefficients and have
the prescribed block degrees. In particular, the block-factorization
theorem applies to real-stable polynomials with fixed block degrees
after a possible overall change of sign.
\end{corollary}

\begin{proof}
Fixed block degrees imply homogeneity. The same-phase theorem of
Y.-B.~Choe, J.~Oxley, A.~Sokal and D.~Wagner
\cite[Theorem~6.1]{ChoeOxleySokalWagner2004} permits multiplication by
a nonzero scalar so that all coefficients are nonnegative. Their theorem
uses the right half-plane; rotating every variable by $i$ gives the
upper-half-plane convention, with only an overall scalar for a homogeneous
polynomial. The resulting polynomial is Lorentzian
\cite[Proposition~2.2]{BrandenHuh2020}, so the theorem gives its factors.
Each factor must be stable: a zero with all its variables in the upper
half-plane would give such a zero of the product.
\end{proof}

\subsection[cellVariableBoundary]{Why we project the cell variables}

For a nonempty word family $\mathcal F$, retain a variable for every cell:
\[
 P_{\mathcal F}(\mathbf z)=\sum_{u\in\mathcal F}\prod_{e\in E}z_{e,u(e)}.
\]
This polynomial has degree one in each row's block of variables. The
theorem shows that it is Lorentzian exactly when $\mathcal F$ is a
Cartesian product of row menus. Indeed, product support is precisely
independent choice of the letters at the rows; conversely such a family
gives a product of nonnegative linear forms, which is even real stable.

If all color quotas are fixed, this can happen only for a single word:
varying one row independently would change the quotas. Thus a nontrivial
fixed-quota family cannot have a Lorentzian full cell polynomial.
The completion-polynomial conjecture instead fixes the hidden cell
weights and retains the positions of one selected color.

\subsection[jointColorPositions]{Which colors can we retain together?}

Let $\mathcal F$ be a nonempty word family with fixed color quotas.
For a set $T$ of colors, keep the variable of every cell whose color lies
in $T$. Give every other cell its own positive weight. We ask when the
resulting position polynomial is Lorentzian for every choice of weights.

Partition the rows into the finest factors on which words can be chosen
independently. For a colored matroid these are the
\href{exchange-graphs.htm#components}{components of the exchange graph
on positions}. A color occurs in a factor if some word uses it there.

\begin{theorem}[Simultaneous color projections]
The joint position polynomial for $T$ is Lorentzian at every positive
hidden weighting if and only if both of the following conditions hold.
\begin{itemize}
\item Each color in $T$ has a Lorentzian position polynomial on its own
at every positive hidden weighting.
\item Each row factor contains at most one color from $T$.
\end{itemize}
The same equivalence holds for real stability.
\end{theorem}

\begin{proof}
Write $P=P_{\mathcal F}$ for the full cell polynomial. Its irreducible
factors over the rationals give exactly the row factors. Indeed, the
least and greatest degrees in each row add under multiplication. Since
$P$ has degree exactly one in each row, that row belongs to just one
irreducible factor. The supports on disjoint rows multiply freely.
Conversely, independent choices on row sets give a product of their
word polynomials, since each word has coefficient one.

Suppose the joint projection has the stated property. Its degrees in
the retained color blocks are fixed. Block factorization therefore gives
\[
 P P_{xy}=P_xP_y
\]
for cells $x,y$ of distinct retained colors, at every positive hidden
weighting. Subscripts denote partial derivatives. This identity holds
in the full polynomial ring since the hidden weights vary independently
over an open set. If an irreducible factor $p$ used both cells, the other
factors would use neither. Cancellation gives $p p_{xy}=p_xp_y$.
But $p$ is prime and both nonzero derivatives have smaller total degree.
This proves the second condition. Evaluating the other retained color
factors at positive weights leaves a positive multiple of the remaining
color factor, which proves the first.

Conversely, consider a row factor containing a retained color $c$.
It has a fixed $c$-quota because the row factors recombine freely.
In the global one-color projection for $c$, differentiate by a supported
monomial from each of the other row factors. Their fixed degrees leave a
positive scalar times the projection on the chosen factor. Derivatives
preserve the stated property. Taking the product over all row factors
now gives the joint projection. This argument applies to both properties.
\end{proof}

\begin{corollary}[Permutation colored matroids]
When each color occurs once, draw the bipartite graph of occurring
row--color cells. A set of colors has a Lorentzian joint position
polynomial for every positive hidden weighting exactly when its colors
lie in distinct connected components. The same sets give real-stable
joint polynomials. In particular, the largest such set has one color
per component.
\end{corollary}

\begin{proof}
Each individual projection is a positive linear form on the possible
rows of its color. These rows are the nonloops of a rank-one matroid,
so they form one matroid component. Joining them for every color gives
exactly the row partition of the bipartite graph. By the
\href{exchange-graphs.htm#components}{component decomposition}, this
is the partition into row factors. The theorem applies.
\end{proof}

In particular, if a Bruhat interval consists of all perfect matchings
on its occurring board, we can read off its jointly retainable colors
from that board. Indeed, exchanging a relative cycle between two
perfect matchings gives two more perfect matchings on the same board.

\begin{example}[Disjoint menus can still interact]
Take all permutations $u$ on four rows with $|u(i)-i|\leq1$. Switching
a relative cycle preserves the board, so this family satisfies symmetric
cycle exchange. Color $1$ can use rows $1,2$, and color $4$ can use rows
$4,3$. In those orders, their joint coefficient matrix at unit weights is
\[
 \begin{pmatrix}2&1\\1&1\end{pmatrix}.
\]
When both colors stay home, rows $2$ and $3$ may either stay home or
exchange places. Each other pair of choices has a unique completion.
The determinant is one. A Lorentzian polynomial of bidegree $(1,1)$
would factor into two linear forms and have coefficient matrix of rank
one. Here the hidden colors couple the two visible choices.
\end{example}

\begin{remark}[Proof status]
The joint-color criterion and its permutation corollary have written
proofs. Independent review and formalization are pending.
\end{remark}
